\documentclass[10pt,a4paper]{amsart}
\usepackage{amsmath,amssymb,amsthm,graphicx}
% For commutative diagrams we use Xy-pic
\usepackage[all]{xy}

% We use 2cell for 2-commutative diagrams.
\xyoption{2cell}
% Theorem environments.
%
\theoremstyle{plain}
\newtheorem{theorem}[subsection]{Theorem}
\newtheorem{proposition}[subsection]{Proposition}
\newtheorem{lemma}[subsection]{Lemma}

\theoremstyle{definition}
\newtheorem{definition}[subsection]{Definition}
\newtheorem{example}[subsection]{Example}
\newtheorem{exercise}[subsection]{Exercise}
\newtheorem{situation}[subsection]{Situation}

\theoremstyle{remark}
\newtheorem{remark}[subsection]{Remark}
\newtheorem{remarks}[subsection]{Remarks}

\numberwithin{equation}{subsection}

% Macros
%
\def\lim{\mathop{\mathrm{lim}}\nolimits}
\def\colim{\mathop{\mathrm{colim}}\nolimits}
\def\Spec{\mathop{\mathrm{Spec}}}
\def\Hom{\mathop{\mathrm{Hom}}\nolimits}
\def\Ext{\mathop{\mathrm{Ext}}\nolimits}
\def\SheafHom{\mathop{\mathcal{H}\!\mathit{om}}\nolimits}
\def\SheafExt{\mathop{\mathcal{E}\!\mathit{xt}}\nolimits}
\def\Sch{\mathit{Sch}}
\def\Mor{\mathop{\mathrm{Mor}}\nolimits}
\def\Ob{\mathop{\mathrm{Ob}}\nolimits}
\def\Sh{\mathop{\mathit{Sh}}\nolimits}
\def\NL{\mathop{N\!L}\nolimits}
\def\CH{\mathop{\mathrm{CH}}\nolimits}
\def\proetale{{pro\text{-}\acute{e}tale}}
\def\etale{{\acute{e}tale}}
\def\QCoh{\mathit{QCoh}}
\def\Ker{\mathop{\mathrm{Ker}}}
\def\Im{\mathop{\mathrm{Im}}}
\def\Coker{\mathop{\mathrm{Coker}}}
\def\Coim{\mathop{\mathrm{Coim}}}

% Boxtimes
%
\DeclareMathSymbol{\boxtimes}{\mathbin}{AMSa}{"02}

%
% Macros for moduli stacks/spaces
%
\def\QCohstack{\mathcal{QC}\!\mathit{oh}}
\def\Cohstack{\mathcal{C}\!\mathit{oh}}
\def\Spacesstack{\mathcal{S}\!\mathit{paces}}
\def\Quotfunctor{\mathrm{Quot}}
\def\Hilbfunctor{\mathrm{Hilb}}
\def\Curvesstack{\mathcal{C}\!\mathit{urves}}
\def\Polarizedstack{\mathcal{P}\!\mathit{olarized}}
\def\Complexesstack{\mathcal{C}\!\mathit{omplexes}}
% \Pic is the operator that assigns to X its picard group, usage \Pic(X)
% \Picardstack_{X/B} denotes the Picard stack of X over B
% \Picardfunctor_{X/B} denotes the Picard functor of X over B
\def\Pic{\mathop{\mathrm{Pic}}\nolimits}
\def\Picardstack{\mathcal{P}\!\mathit{ic}}
\def\Picardfunctor{\mathrm{Pic}}
\def\Deformationcategory{\mathcal{D}\!\mathit{ef}}
\title{The multiplicity of a two-periodic complex built from a nilpotent endomorphism vanishes}
\author{The Stacks Project Authors}
\date{stacks-project:chow.tex:local-snapshot-20261011; GNU FDL 1.2 or later}
\begin{document}
\maketitle
\noindent\textit{Editorial note.} All delivered mathematical text is the author's own, copied from the pinned Stacks Project chapter file chow.tex; no statement or proof step was rewritten, reordered or removed.\par\medskip
\noindent\textit{Editorial note.} The preamble is the chapter's own native preamble excerpt used for this source: the Xy-pic setup at lines 14-18 and the theorem declarations with the mathematical macros at lines 154-217 of sources/stacks-chow/preamble.tex. The quoted class line and the xr-hyper externaldocument block are not included because the assembler supplies a plain amsart wrapper; no macro was invented or redefined.\par\medskip
\noindent\textit{Editorial note.} The counted statement and proof contain no cross reference at all, so the assembled document defines every label it uses and no external reference stub is needed.\par\medskip
\noindent\textit{Allowed standard prerequisites (editorial).}\begin{itemize}
\item Two-periodic and (2,1)-periodic complexes and their multiplicity e\_R. Same chapter, Definition environment labelled definition-periodic-complex (source line 110) together with the displayed definition of e\_R at source lines 200-205 (label equation-multiplicity-coker-ker, line 203). Supplies the objects M, N, phi, psi, the notation e\_R(M, phi\textasciicircum{}i, phi\textasciicircum{}\{n-i\}) and the convention phi\textasciicircum{}0 = id\_M used in the first line of the counted proof; the vanishing argument is the counted content.
\item Length is additive in short exact sequences. Algebra, Lemma labelled algebra-lemma-length-additive (canonical Stacks tag 00IV). Not needed by this card, which works with lengths of the explicit modules K\_i/t\textasciicircum{}j M directly; recorded because the sibling lemma of the same section uses it.
\end{itemize}\medskip
\begin{lemma}
\label{lemma-powers-period-length-zero}
Let $R$ be a ring. Let $M$ be an $R$-module.
Let $\varphi : M \to M$ be an endomorphism and $n > 0$
such that $\varphi^n = 0$ and such that $\Ker(\varphi)/\Im(\varphi^{n - 1})$
has finite length as an $R$-module.
Then
$$
e_R(M, \varphi^i, \varphi^{n - i}) = 0
$$
for $i = 0, \ldots, n$.
\end{lemma}

\begin{proof}
The cases $i = 0, n$ are trivial as $\varphi^0 = \text{id}_M$ by convention.
Let us think of $M$ as an $R[t]$-module where multiplication by $t$
is given by $\varphi$. Let us write
$K_i = \Ker(t^i : M \to M)$ and
$$
a_i = \text{length}_R(K_i/t^{n - i}M),\quad
b_i = \text{length}_R(K_i/tK_{i + 1}),\quad
c_i = \text{length}_R(K_1/t^iK_{i + 1})
$$
Boundary values are $a_0 = a_n = b_0 = c_0 = 0$.
The $c_i$ are integers for $i < n$ as $K_1/t^iK_{i + 1}$
is a quotient of $K_1/t^{n - 1}M$ which is assumed to have finite length.
We will use frequently that $K_i \cap t^jM = t^jK_{i + j}$.
For $0 < i < n - 1$ we have an exact sequence
$$
0 \to
K_1/t^{n - i - 1}K_{n - i} \to
K_{i + 1}/t^{n - i - 1}M \xrightarrow{t} K_i/t^{n - i}M
\to K_i/tK_{i + 1} \to 0
$$
By induction on $i$ we conclude that $a_i$ and $b_i$ are
integers for $i < n$ and that
$$
c_{n - i - 1} - a_{i + 1} + a_i - b_i = 0
$$
For $0 < i < n - 1$ there is a short exact sequence
$$
0 \to
K_i/tK_{i + 1} \to
K_{i + 1}/tK_{i + 2} \xrightarrow{t^i}
K_1/t^{i + 1}K_{i + 2} \to
K_1/t^iK_{i + 1} \to 0
$$
which gives
$$
b_i - b_{i + 1} + c_{i + 1} - c_i = 0
$$
Since $b_0 = c_0$ we conclude that $b_i = c_i$ for $i < n$.
Then we see that
$$
a_2 = a_1 + b_{n - 2} - b_1,\quad
a_3 = a_2 + b_{n - 3} - b_2,\quad \ldots
$$
It is straightforward to see that this implies $a_i = a_{n - i}$ as desired.
\end{proof}

\end{document}
